% GENERATED FILE. Edit canonical lessons, then run npm run generate:books.
% canonical-source-hash: fnv1a64:61b3d23c489e58bb
% canonical-lessons: RU-C151-mesto, RU-C151-tsentr, RU-C151-okraina, RU-C151-rayon, RU-C151-prospekt

\chapter{Place, Centre, Outskirts, District, Avenue}
\label{ch:ru-place-centre-outskirts-district-avenue}
\hlchaptermodality{151}

\begin{chapteropening}
\textbf{By the end of this chapter:} \emph{I can say a place, the centre, the outskirts, a district and an avenue.}

Complete the last lesson of chapter 151: I can say a place, the centre, the outskirts, a district and an avenue.
\end{chapteropening}

\section[\texorpdfstring{\ru{место} (mésto)}{место (mésto)}]{\ru{место} (mésto) — a place; a seat}

\label{lesson:RU-C151-mesto}



\begin{quote}
Before the new one: say the Russian for at last, then the Russian for still; more.
\end{quote}

\subsection*{You'll want to know: \ru{место}}
\textbf{\ru{место}} (\emph{mésto}) — \textquotedblleft{}a place; a seat\textquotedblright{}.

\begin{cousinweb}
\textbf{\ru{место}} $\leftarrow$ Proto-Slavic \emph{*město}, \textquotedblleft{}a place\textquotedblright{} — inherited. Its sense \textquotedblleft{}town\textquotedblright{}, in Ukrainian \emph{misto} and Polish \emph{miasto}, Vasmer explains as a copy of Middle High German \emph{stat}, \textquotedblleft{}place; town\textquotedblright{}.
\end{cousinweb}

\begin{grammarlens}[title={What you've built}]
One more word for this chapter.
\end{grammarlens}

\subsection*{Your turn}
\begin{itemize}
\raggedright
  \item \emph{Say it:} \emph{mésto}
  \item \emph{Say it:} \emph{mésto}, once more
  \item \emph{Say it:} say \emph{mésto}
  \item \emph{Recall:} say the Russian for at last, then the Russian for still; more, then say \emph{mésto} again
\end{itemize}

\begin{tcolorbox}[breakable,colback=teal!4,colframe=teal!35!black,title={Before you move on}]
What does \ru{место} mean? (\textquotedblleft{}A place; a seat\textquotedblright{}.) Say it once more.
\end{tcolorbox}
\section[\texorpdfstring{\ru{центр} (tsentr)}{центр (tsentr)}]{\ru{центр} (tsentr) — the centre}

\label{lesson:RU-C151-tsentr}



\begin{quote}
Before the new one: say the Russian for still; more, then the Russian for a place; a seat.
\end{quote}

\subsection*{You'll want to know: \ru{центр}}
\textbf{\ru{центр}} (\emph{tsentr}) — \textquotedblleft{}the centre\textquotedblright{}.

\begin{cousinweb}
\textbf{\ru{центр}} — recorded from the 1600s; it came through German \emph{Zentrum} from Latin \emph{centrum}, Greek \emph{kéntron}, \textquotedblleft{}a sharp point; the fixed leg of a compass\textquotedblright{} (Vasmer).
\end{cousinweb}

\begin{grammarlens}[title={What you've built}]
One more word for this chapter.
\end{grammarlens}

\subsection*{Your turn}
\begin{itemize}
\raggedright
  \item \emph{Say it:} \emph{tsentr}
  \item \emph{Say it:} \emph{tsentr}, once more
  \item \emph{Say it:} say \emph{tsentr}
  \item \emph{Recall:} say the Russian for still; more, then the Russian for a place; a seat, then say \emph{tsentr} again
\end{itemize}

\begin{tcolorbox}[breakable,colback=teal!4,colframe=teal!35!black,title={Before you move on}]
What does \ru{центр} mean? (\textquotedblleft{}The centre\textquotedblright{}.) Say it once more.
\end{tcolorbox}
\section[\texorpdfstring{\ru{окраина} (okráina)}{окраина (okráina)}]{\ru{окраина} (okráina) — the outskirts}

\label{lesson:RU-C151-okraina}



\begin{quote}
Before the new one: say the Russian for a place; a seat, then the Russian for the centre.
\end{quote}

\subsection*{You'll want to know: \ru{окраина}}
\textbf{\ru{окраина}} (\emph{okráina}) — \textquotedblleft{}the outskirts\textquotedblright{}.

\begin{cousinweb}
\textbf{\ru{окраина}} is \emph{\ru{о}-}, \textquotedblleft{}around\textquotedblright{}, with \emph{\ru{край}}, \textquotedblleft{}edge\textquotedblright{}, and \emph{-\ru{ина}}: the land at the edge. Vasmer connects \emph{\ru{край}} with \emph{\ru{кроить}}, \textquotedblleft{}to cut\textquotedblright{}: an edge is where something is cut off.
\end{cousinweb}

\begin{grammarlens}[title={What you've built}]
One more word for this chapter.
\end{grammarlens}

\subsection*{Your turn}
\begin{itemize}
\raggedright
  \item \emph{Say it:} \emph{okráina}
  \item \emph{Say it:} \emph{okráina}, once more
  \item \emph{Say it:} say \emph{okráina}
  \item \emph{Recall:} say the Russian for a place; a seat, then the Russian for the centre, then say \emph{okráina} again
\end{itemize}

\begin{tcolorbox}[breakable,colback=teal!4,colframe=teal!35!black,title={Before you move on}]
What does \ru{окраина} mean? (\textquotedblleft{}The outskirts\textquotedblright{}.) Say it once more.
\end{tcolorbox}
\section[\texorpdfstring{\ru{район} (rayón)}{район (rayón)}]{\ru{район} (rayón) — a district}

\label{lesson:RU-C151-rayon}



\begin{quote}
Before the new one: say the Russian for the centre, then the Russian for the outskirts.
\end{quote}

\subsection*{You'll want to know: \ru{район}}
\textbf{\ru{район}} (\emph{rayón}) — \textquotedblleft{}a district\textquotedblright{}.

\begin{cousinweb}
\textbf{\ru{район}} comes from French \emph{rayon}, \textquotedblleft{}a ray, a sector\textquotedblright{}, from Latin \emph{radius}, \textquotedblleft{}a spoke, a ray\textquotedblright{} (Vasmer).
\end{cousinweb}

\begin{grammarlens}[title={What you've built}]
One more word for this chapter.
\end{grammarlens}

\subsection*{Your turn}
\begin{itemize}
\raggedright
  \item \emph{Say it:} \emph{rayón}
  \item \emph{Say it:} \emph{rayón}, once more
  \item \emph{Say it:} say \emph{rayón}
  \item \emph{Recall:} say the Russian for the centre, then the Russian for the outskirts, then say \emph{rayón} again
\end{itemize}

\begin{tcolorbox}[breakable,colback=teal!4,colframe=teal!35!black,title={Before you move on}]
What does \ru{район} mean? (\textquotedblleft{}A district\textquotedblright{}.) Say it once more.
\end{tcolorbox}
\section[\texorpdfstring{\ru{проспект} (prospékt)}{проспект (prospékt)}]{\ru{проспект} (prospékt) — an avenue}

\label{lesson:RU-C151-prospekt}



\begin{quote}
Before the new one: say the Russian for the outskirts, then the Russian for a district.
\end{quote}

\subsection*{You'll want to know: \ru{проспект}}
\textbf{\ru{проспект}} (\emph{prospékt}) — \textquotedblleft{}an avenue\textquotedblright{}.

\begin{cousinweb}
\textbf{\ru{проспект}} probably came through German \emph{Prospekt}, \textquotedblleft{}a view, a vista\textquotedblright{}, from Latin \emph{prōspectus}, \textquotedblleft{}a view\textquotedblright{} (Vasmer): a long straight avenue gives a view; \emph{\ru{прешпект}} was the folk form.
\end{cousinweb}

\begin{grammarlens}[title={What you've built}]
One more word for this chapter.
\end{grammarlens}

\subsection*{Your turn}
\begin{itemize}
\raggedright
  \item \emph{Say it:} \emph{prospékt}
  \item \emph{Say it:} \emph{prospékt}, once more
  \item \emph{Say it:} say \emph{prospékt}
  \item \emph{Recall:} say the Russian for the outskirts, then the Russian for a district, then say \emph{prospékt} again
\end{itemize}

\begin{tcolorbox}[breakable,colback=teal!4,colframe=teal!35!black,title={Before you move on}]
What does \ru{проспект} mean? (\textquotedblleft{}An avenue\textquotedblright{}.) Say it once more.
\end{tcolorbox}
